\section{Introduction}\label{sec:Intro}
This paper presents an extension of the work by \cite{OellerichBuescherDegel2023_1000162443}, which established the foundation of the generalized isocontouring approach. While \cite{OellerichBuescherDegel2023_1000162443} employed an arbitrary plane in a three-dimensional space, the present work considers a hyperplane that is normal to one of the coordinate axes. Although this choice reduces the level of generality, it introduces a systematic and scalable method for extending the approach to spaces of arbitrary dimensionality. In practice, isocontouring is most often required for hyperplanes normal to coordinate axes, so this restriction has limited impact on typical use cases. Moreover, this modification significantly enhances computational efficiency and improves the robustness of subsequent calculations based on the isocontouring results, as the resulting representations are unique and thereby avoid, for example, unintended weighted Euclidean averaging.